% ============================================================
%  Hans Brøchner, "Recollections of Søren Kierkegaard"
%  [Erindringer om Søren Kierkegaard]
%  Det nittende Aarhundrede. Maanedsskrift for Literatur og Kritik,
%  ed. Georg and Edvard Brandes, vol. 5, March 1877, printed pp. 337--374.
%  Written 27 Dec. 1871 -- 10 Jan. 1872; published after Brøchner's death.
%  TRANSLATION (English).
%  Translated from transcription.tex (Danish) in this directory, whose
%  header records how its text was established (every page read twice at
%  the image).  Method: ../../../TRANSLATION-PLAYBOOK.md.  No published
%  translation or later edition was consulted (T. H. Croxall's Glimpses and
%  Impressions of Kierkegaard, 1959, and Steen Johansen's 1953 edition are
%  in copyright and were not opened).
%
%  ---- REGISTER ----
%  Readable modern English, scholarly, moderately literal; American
%  spelling.  The piece is a memoir -- anecdote, conversation, character
%  sketch -- and is rendered as one: Brøchner's affection and irony, and
%  the point of Kierkegaard's reported sayings, are kept.  Long periods are
%  broken only where English will not carry them, never across a page
%  marker or a footnote anchor.
%
%  ---- TERMINOLOGY ----
%  Aligned with the other Brøchner, Brandes and Nielsen translations in
%  this repo:
%    Tro / Viden = faith / knowledge (one ordinary "belief", p. 361, is
%      commented at the site);  Erkendelse = cognition;  Videnskab =
%      science (Wissenschaft);  Aand = spirit;  Kristendom = Christianity;
%    Opbyggelse = upbuilding;  den Enkelte = the single individual;
%    Forfattervirksomhed = authorship;  det Bestaaende = the established
%      order;  Præst = pastor;  Livsanskuelse = view of life.
%  Kierkegaard's works are given in their standard English titles (Either/Or,
%  Concluding Postscript, Works of Love, Upbuilding Discourses in Various
%  Spirits, Practice in Christianity, The Moment; Corsaren = The Corsair),
%  as in this repo's translation of Brøchner's 1855 article; Διαψάλματα
%  stays in Greek.  All other Danish titles (Fædrelandet, Genboerne, ...)
%  stay in Danish.  Names unchanged; "S. K." and "K." as Brøchner writes
%  them; Hr. / Fru = Mr. / Mrs.; Danish titles of rank (Etatsraad,
%  Kammerraad, ...) untranslated.
%
%  ---- CONVENTIONS ----
%  - Page markers \opage{N} and "% --- p. N ---" at the same joints as in
%    transcription.tex, so the two can be read side by side; a word the
%    Danish divides across a page turn is kept whole in English, with the
%    marker before or after it.
%  - The same paragraphs, numbered sections (\secnum), \emph spans
%    (letterspacing in the original), footnote and centre blocks as the
%    Danish.  The duplicated section number "3." (p. 340) is kept as
%    printed; section numbers are structure, not sense.
%  - Quotation marks „...“ -> ``...''.  Latin, German and French stay in
%    their language; Greek verbatim.
%  - Omission lines and the points or asterisks standing for withheld names
%    are reproduced as printed.
%
%  ---- DEFECTS: DANISH AS PRINTED, ENGLISH CORRECT ----
%  The transcription is diplomatic and carries every printer's error (logged
%  in its header).  The translation renders the SENSE and logs the
%  divergence in a % comment at the site; where the print omits an opening
%  or closing quotation mark (pp. 348, 350, 352, 374) the English supplies
%  it, so the English quotation marks balance (74 / 74).
%
%  ---- HOW THE TRANSLATION WAS CHECKED (2026-09-23) ----
%  (1) Four translators, one per ten-page batch, each re-reading its English
%  against the Danish sentence by sentence for omissions, negations, modals
%  and tone before returning.  (2) Mechanical audit (tr_audit.py): per
%  printed page, the same \opage joints, \emph/\textit/\textbf spans,
%  footnote, section numbers, centre blocks and paragraph breaks as
%  transcription.tex, and an English/Danish word ratio within 0.95--1.60:
%  38 pages, 0 flagged (ratios 1.04--1.19).  joints.py: both batch seams
%  (pp. 347, 367) run on mid-sentence without a spurious paragraph break.
%  (3) Independent review: four readers who had not translated the pages
%  read every sentence of the Danish against the English.  No omission
%  found; 9 corrections applied -- p. 345 "paa Glatis" (tripping up, not an
%  outing on the ice), "hverken vidste ud eller ind" (a calque replaced),
%  "Strasburg" as printed; p. 357 "Antiker" = antique sculptures (false
%  friend); p. 363 "den store Skikkelse" (Poul Møller's physical bulk),
%  "Fru" = Mrs.; p. 364 "i den" refers to grief, not work; p. 371 it is
%  Brøchner who walks past Mynster; p. 374 "skulde have talt" -- the
%  conversation just ended was the last.
%
% ============================================================
\documentclass[12pt,a4paper,openany]{book}
\usepackage[utf8]{inputenc}
\usepackage[T1]{fontenc}
\usepackage{libertinus}
\usepackage{libertinust1math}
\usepackage{textalpha}   % Greek: Διαψάλματα (p. 348), μεσότης (p. 361)
\usepackage{graphicx}
\usepackage[english]{babel}
\usepackage[protrusion=true,expansion=true]{microtype}
\usepackage{setspace}
\usepackage[margin=1.2in]{geometry}
\usepackage{enumitem}
\usepackage{fancyhdr}
\setlength{\headheight}{28pt}
\addtolength{\topmargin}{-13.5pt}

% Footnotes as the 1877 printing sets them: *).
\renewcommand{\thefootnote}{*)}

% The article's numbered sections: the number alone, centred on its line,
% exactly as in transcription.tex.
\newcommand{\secnum}[1]{\par\begin{center}#1\end{center}\par\nobreak}

\usepackage{hyperref}
\hypersetup{
  pdftitle={Recollections of Søren Kierkegaard},
  pdfauthor={Hans Brøchner},
  hidelinks
}

\pagestyle{fancy}
\fancyhf{}
\fancyhead[LE,RO]{\thepage}
\fancyhead[RE]{\textit{H. Brøchner: Recollections of Søren Kierkegaard}}
\fancyhead[LO]{\textit{Det nittende Aarhundrede, March 1877}}

\setstretch{1.3}

\title{\textbf{Recollections of Søren Kierkegaard}\\[8pt]
  \large [Erindringer om Søren Kierkegaard]}
\author{Hans Brøchner}
\date{\textit{Det nittende Aarhundrede. Maanedsskrift for Literatur og
  Kritik}, vol.~5, March 1877, pp.~337--374\\[10pt]
  \small Translated from the Danish}

% ---- original-page markers ----
% \opage{N} marks where printed page N begins, at the same joint as in
% transcription.tex (mid-sentence where the page turns mid-sentence).
\usepackage{marginnote}
\newcommand{\opage}[1]{\marginnote{\footnotesize\textit{[#1]}}}
\newcommand{\apage}[1]{\marginnote{\footnotesize\textit{[#1]}}}

\begin{document}
\maketitle
\thispagestyle{empty}
\clearpage

% --- p. 337 ---
\begin{center}\textbf{\large \opage{337}Recollections of Søren Kierkegaard.}\\[4pt]\rule{2cm}{0.4pt}\end{center}

I have wanted to write down some of the recollections I
have of Søren Kierkegaard from a long run of years, and
which might be of some interest to others as well. They
might perhaps one day, when S. K. finds a competent biographer,
help to make his portrait in certain respects
more finished and more alive. I write down my recollections
precisely as material for a portrait of K., and
therefore without any attempt, at this stage, to work them
up into a whole or to give them a rounded finish by
gathering them into groups. They are given just as
they have been preserved, after the first impression, as \emph{separate}
memories in a faithful memory.

\secnum{1.}

% Danish prints `K.‘s' with a turned comma for the apostrophe; rendered as the sense.
My first meeting with S. K. falls in the very earliest
days after I had become a student (1836). I saw
him then in the house of my old uncle, the wholesale merchant M. A. Kierkegaard,
at a party given on the occasion of the present
Bishop Kierkegaard's engagement (or wedding?) to his
first wife. I saw S. K. there without knowing who
he was; I had only been told that he was Dr.
K.'s brother. He spoke only very little that evening; he
kept, I believe, essentially to observing. I received only a
% --- p. 338 ---
\opage{338}definite impression of his exterior, which rather
amused me. He was then 23 years old, had
something very irregular about his whole figure and a
remarkable coiffure. His hair rose in a tousled crest
% `et Kvarter': a quarter of an ell (c. 16 cm); rendered `six inches'.
nearly six inches above his forehead, and it gave him
a strangely wild look. I got the idea, without my
knowing how, that he was a merchant's clerk
--- perhaps because the family was a family of merchants ---
and under the impression of his odd
exterior I at once completed it: he must be in a draper's shop.
I have often laughed heartily since at my own acumen.

\secnum{2.}

Then I remember, in the time just after, now and then
having seen S. K. at the same place. I now came to
talk with him, and soon noticed that he belonged somewhere
other than behind a counter with the ell-measure in his hand.
Amid the very plain surroundings there he amused himself
by putting forward little paradoxes. Thus I remember
that he once astonished his relatives by declaring, with
the most serious face in the world, that he found
the old A. B. C. we used in our childhood to be one
of the most interesting books there were, and that he
still read in it quite regularly, and got much profit
from it. --- One evening he joined in a game of boston. It was played
there with variations, one of which consisted in this, that
3 of the players were dealt 17 cards each, the
fourth, on the other hand, only one, and that the 3 others then each gave up
4 cards to the last. In that house, where card-playing was taken
with devout seriousness, the situation of the 4th man in this game
was always regarded as most lamentable,
and S. K. again caused great astonishment by expounding that
it was the most excellent situation that could
be imagined, and that his only wish was always to be able to be
the fourth man so situated. That would be the
most piquant thing of all.

% --- p. 339 ---
\secnum{3.}

\opage{339}In 1837 I met K. now and then at a restaurant.
He no longer lived with his father at that time, but in a
house in Løvstræde, I believe the one where Reitzel's
bookshop now (1871) is. He usually took his supper at a restaurant,
and I can remember that I, when I now and
then treated myself to ``a half boeuf,'' was astonished
at the luxury S. K. displayed when he took his
supper with a half bottle of wine and the like. We now talked
together more frequently, and he showed himself very kind toward
me. One evening he asked me what I was reading in the way of
aesthetic writings. On that occasion I came to
mention that I had access to only a little, and that I
was quite unacquainted with many fields. He asked
me whether I knew the writings of the German Romantics,
and I had to say that I did not. He then urged me
to come home with him and lent me a book by
Eichendorff: Dichter und ihre Gesellen, --- a book I
later bought at the auction of his books in remembrance
of this meeting with him. I then remember that when,
after a fortnight or so, I brought it back to him, and
was just about to make an excuse because I, who at that time
was studying theology with passion, had kept it
so long, I was disconcerted when he received me
with the question whether I had \emph{already} read it.
He had evidently, with his sharp eye, been able to see that
the opposite word was hovering on my tongue, and he
amused himself by embarrassing me in a good-natured way;
which at that time was exceedingly easy. I
blushed so easily, and that sight always appealed to him
in ``the young person.'' --- From my visit to K., when
he lent me the book, I remember two things. One is
my astonishment at his large library, which quite
impressed me. The other was a peculiarity of K.'s. He
went out again when we had fetched the book, and, as
we were leaving, blew out a wax candle he had lit. He
% --- p. 340 ---
\opage{340}explained to me then that he always did it with a certain
caution and at some distance, because he had
reflected himself into the notion that the smoke of the candle
was dangerous to inhale and could harm his
chest.

% Danish prints a second section `3.' here (it should be 4; no 4 is printed); the number is reproduced as printed.
\secnum{3.}

In 1838 S. K.'s father died. Him I saw only
once, at my old uncle Kierkegaard's. He was
a considerable personality, whose outward appearance imprinted itself vividly
on my memory. He was walking up and down
the floor of my uncle's dining room with him, in conversation. He
walked somewhat stooped; his bearing matched the thoughtful,
serious expression of his face; the gray hair was brushed
back behind the ear, so that the form of the face stood out sharply.
--- The reverence and love S. K. bore his
father have found their living expression in his writings.
--- The old man was a very punctual man, exact in money matters,
thrifty in everyday things, but openhanded when occasion
offered. S. K. once told me how
he, as a young student, had made a journey in Nordsjælland,
and how the old man had fitted him out
generously with travel money, and then had even surprised him
by sending him, at one of the last stages of the journey, a
registered letter with 50 rix-dollars. Of his father's punctuality
S. K. once, in a conversation with Mag. Adler, used the expression:
% `født i Terminen': Termin is both the quarter day on which rents and payments fell due and the expected date of a birth; the English keeps only the first sense.
``my father was born on quarter day.'' --- When S. K. left
his father's house, the old man gave him yearly a sum that
was very ample by the standards of the time; if I am not
much mistaken, it was 800 rix-dollars a year.

\secnum{5.}

My late sister Hansine was in town in 1839 as a young girl.
Once, when I happened to mention S. K.,
she told me that she had seen him when, before her
confirmation, she was here in town on a visit to old Uncle
% --- p. 341 ---
\opage{341}Kierkegaard. Our cousins had one day, when S. K.
was to come there, given her the advice not to have anything
to do with him, since he was ``a dreadfully spoiled
and naughty boy, who always clung to his mother's skirts.''
My sister was then 13 years old, S. K. 15.
Later she read with great pleasure some of his upbuilding
writings, especially ``Works of Love,'' and
could then sometimes, with the smile peculiar to her,
remind me of the cousins' admonition to her.

\secnum{6.}

Shortly after S. K. had taken his theological degree
in 1840, I talked with him. He told me that
his father had constantly wished him to take the
theological examination, and that they had very often discussed
the matter. ``As long as Father was alive, I could
nevertheless defend my theory that I ought not to take
it, but when he was dead, and I also had to take over
his \textit{partes} in the debate, then I could no longer hold
my own, and so I had to make up my mind to read for the examination.''
This S. K. had then done with great energy. He had
taken ``Sin Brøchner'' as his coach, had worked his way
through the driest disciplines, made excerpts from church history,
learned ``the lists of popes'' by heart and the like. His coach
was very pleased with him. Shortly after S. K.
had taken the examination, Peter Stilling (later Dr. phil.)
came to Sin Brøchner and wanted to be coached. He remarked that
he thought he could be finished in a year and a half;
``S. K. hadn't read any longer than that, after all.'' ``Oh, to be sure,'' said
old Br., who did not excel in politeness, ``don't you go
imagining that! with S. K. it was another matter, he
% Danish prints a semicolon after `erindrer' where a comma is wanted; rendered as the sense.
could do anything!'' --- I remember that in 1839 I attended
theological writing exercises together with S. K. under Clausen.
% Danish prints `eiler' for `eller'; rendered as the sense (`or').
He took part in them, however, only for the first 2 or 3 sessions.
It was told that he had once, several years earlier,
taken part in Clausen's writing exercises and once, instead of
% --- p. 342 ---
\opage{342}answering the set question, had analyzed it and shown that
there was no sense to be found in it. At this Cl. had been
offended, and S. K. had stopped going to the writing exercises.
--- At the written theological examination S. K. was placed No.
4 in the grading. Ahead of him stood M. Wad, Warburg
and Chr. F. Christens. The examiners, however, declared
to these that K.'s papers bore witness to a far greater
maturity and development of thought than any of the others';
but theirs contained a richer measure of positive theological matter.

\secnum{7.}

% Danish prints `i hvilke' (plural) with the singular `Samtale'; rendered as the sense (`in which').
In the same conversation in which S. K. told me
how he had come to take the theological examination,
he also spoke of his father's remarkable calm and objectivity.
Thus the old man had once said to S. K.:
``it would really be a stroke of luck for you if I were dead;
then perhaps you might still amount to something; you will
not, as long as I live.'' --- (Another trait of his
father was told me by my late cousin Peter Kierkegaard.
When S. K.'s mother died, his father was deeply
grieved. At the funeral Bishop Mynster came up
to him and, moved, expressed his sympathy. The old man
listened to Mynster, whom he esteemed highly, without betraying his grief,
and when Mynster fell silent, he said: ``Shall we
not then, Your Right Reverence, go into the room there
and drink a glass of wine.'' Mynster, who knew him, did not
misunderstand this remark.)

\secnum{8.}

It was while S. K. was in Berlin in 1841--42 that I
presented myself for the theological examination, to which, on account
of my unorthodox views, I was refused admission.
I was corresponding at that time with Chr. F. Christens, who
was also in Berlin, and through him I sent and received
greetings, also to and from Kierkegaard. Shortly after
his return home I met him in the street. He was
% --- p. 343 ---
\opage{343}very friendly toward me, talked about my situation and my
studies, and told me, which pleased me, that in
Berlin he had longed for several things back home, among them
also to see me. Such a thing he could say with
a peculiarly winning expression. His smile and his
glance were indescribably expressive. He had a way of his own
of greeting at a distance with a glance. It was only
a small movement of the eye, and yet it expressed so much.
There could be something infinitely gentle and loving in his
glance, but also something provoking and teasing. He could,
with one glance at a passer-by, irresistibly ``put
himself in rapport'' with him, as he put it. Whoever
met this glance was either attracted or repelled,
made embarrassed, uncertain or irritated. I have walked the whole
length of a street with him while he expounded to me
how, by thus putting oneself in rapport with
passers-by, one could make psychological studies, and
while he expounded the theory, he put it into
practice with nearly every person we met. There
was not one on whom his glance did not make a visible
impression. On the same occasion he astonished me by
the ease with which he struck up a conversation with
any number of people, in a few exchanges took up an
earlier conversation and carried it a step further to a
point where it could be taken up again on occasion.
The occasion for these experiments was that one day I
had been walking ahead of him, absorbed in thought, so that I
had not heard him calling me, and had not noticed
him tapping me on the shoulder. When I finally noticed
him, he said that one did wrong thus to sink
into oneself and not to make the observations one could
make in such rich measure. Then, to show me the method,
he pulled me along with him up and down several streets, and astonished
me with his talent for psychological experiment.
He was always interesting to walk with, but in one respect
walking with him had its difficulty.
% --- p. 344 ---
\opage{344}Because of the irregularity of his movements, which
was no doubt connected with his crookedness, one could
never keep to a straight line when walking with him;
one was always pushed, little by little, in toward the houses
and the cellar steps, or out toward the gutter. When
he, in addition, gestured with his arm and his cane,
it became even more of a march over obstacles.
One had now and then to seize an opportunity to
get round to his other side in order to gain
the room one needed.

\secnum{9.}

The winter S. K. spent in Berlin, he met there
many Danes, not a few of whom had
come there to hear Schelling. Among them was
my late friend Christian Fenger Christens. Of him
S. K. spoke with great appreciation; he stated later
to me that Christens was the ablest of all the Danes
who had been in Berlin that winter, and among them
were, after all, such men as Lic. jur. (later Minister)
A. F. Krieger, and the judge advocate (now head of department in the Ministry of Church and Education)
Carl Weiss. Christens told me various
little traits of K., of his helplessness when he
had to speak German, especially when he had to speak of things
of daily life (thus on the first evening in Berlin
it was not possible for him to make it understood that he
wanted a candlestick); --- of how his landlord
% `trak ham op ... Optrækkerierne': taken as `overcharged' (trække op = to fleece), which the landlord offsets by promoting K. in title.
overcharged him and, as the overcharging grew, by way of compensation
let him advance from Magister to Doctor and then
to Professor; --- something Kierkegaard himself was amused
at, since he no doubt guessed the reason, --- of his embarrassment
when Weiss, who set store by an elegant dinner, had one day
taken him into a first-class restaurant,
where a number of elegantly dressed gentlemen in black coats,
with white neckcloths and cravats, stood in the corners
of the room, and after K. had respectfully greeted them
% --- p. 345 ---
\opage{345}and sat down at a table with W., came leaping up and turned out
to be waiters etc. --- Christens also told me
how Kierkegaard in Berlin amused himself by
getting Lic. theol. * onto slippery ground. * had already been
abroad for a long time, had heard lectures in
Strasburg and elsewhere and was very self-important on
account of all the wisdom he had gathered. S. K.
seldom failed, when they met, to ask him
to share with him the most essential scientific results
of his journey abroad. When * then formulated some proposition
or other as a result gained, K. always knew how
either to show him that it was something very old,
or, by asking little questions about ``something that was not
quite clear to him'' in what had been put forward, to bring the
poor licentiate, who was only slightly gifted as a thinker,
into such confusion that he did not know which way
to turn. When K. had then got him to that point,
it suddenly occurred to him that he was very
busy and absolutely had to go somewhere or other.
Then he left the licentiate standing in the midst of the confusion, to the great
amusement of the others. --- In K.'s Either/Or
this method is depicted, applied precisely to a licentiate. Might
he not have had his Berlin acquaintance in
mind when he wrote those pages?

\secnum{10.}

Not long after K.'s return home he once talked
with me about my studies. I told him what
Greek authors I was reading at the time, and touched on the
interest that was awakening in me more and more for
Greek poetry and philosophy. He encouraged me to
continue these studies and spoke warmly of the significance
of the Greek for our time. He used the expression:
``Greece ought always to have a representative
accredited to our time, just as one power has a representative
accredited to another to make its interests
% --- p. 346 ---
\opage{346}felt.'' --- K. himself had a good knowledge of Greek and
had read not a little. How finely he apprehended the
Greek spirit, his writings bear witness.

\secnum{11.}

In 1842 a small pseudonymous work appeared: ``Johan
Ludvig Heiberg efter Døden.'' I knew of the work
before it appeared, and mentioned it once to S. K. He
was not pleased that Heiberg was made the object
of such a joke, and spoke with much
warmth of his importance as an aesthetician and especially
as his generation's teacher of aesthetics in our fatherland\footnote{K. placed H. above all the contemporary aestheticians in Germany (NB. Vischer had not yet at that time come forward except with lesser things.)}. Later, after Either/Or had appeared, and
Heiberg had written his well-known review of it,
S. K. once spoke to me about him and made his displeasure
at H.'s conduct plain to see. He
acknowledged H.'s importance as an aesthetician, but now also
strongly stressed his limitations. ``I would be able
to bring forward a whole series of aesthetic problems
of which H. has not an inkling.'' By this he was no doubt alluding
to the problems that have come forward in ``Stages
on Life's Way'' and ``Concluding Postscript.''

\secnum{12.}

In those years I frequently met S. K. when in the evening I
went for a walk. As a rule Frederiksberg
Have was the goal of my walk. K. too went that far,
but only to the entrance of the garden, where the little flower beds
are on either side of the path, narrow at that time,
that leads from the gate to the first open space. There he
breathed in for a few moments the scent of the flowers, and
% --- p. 347 ---
\opage{347}then took the recollection of this ``moment'' away with him. That was
how he liked to conclude and to set limits: his walk had a definite goal,
but this goal was, as it were, only touched; he did not linger over it,
and he took away, as it were, only a motif of enjoyment that could be
worked up ideally.

A small comic recollection is attached to one of these evening walks. As
we were going back from the garden to the avenue and had to pass between
the posts that stand at the end of the path between the trees that is set
aside for pedestrians, the posts had just been painted. K. did not notice
this, but rested his hand on them as he passed between them. He then
noticed that his hand was covered in paint, and now, while carrying on
with his conversation, made a series of unsuccessful attempts, with the
indispensable cane under his arm, to wash the paint off by rubbing the
drooping leaves of the trees, damp with dew, between his hands. His
exertions brought many a smile to the faces of the passers-by.

\secnum{13.}

At that time I also saw S. K. on horseback now and then. He had learned
to ride so as to be able to take this exercise and to make short
excursions independently of coachmen and the like. He cut no very good
figure on a horse. One could see from his posture that he did not credit
himself with any great ability to control it, should it take it into its
head to turn rebellious. He sat stiffly on the horse, and one had the
impression that he was constantly remembering the riding master's
instructions. He can hardly have had much freedom to pursue his thoughts
and his fancies on horseback. And indeed he soon gave up this sport and
preferred to drive when he wanted to visit his favorite spots in the woods
around Copenhagen. In the years when his authorship was at its most
intense, these outings were one of the means he used to keep
% --- p. 348 ---
\opage{348}himself fresh, and to bring himself into the mood that the
production demanded.

\secnum{14.}

One single expression of K.'s from that time I cite as characteristic. He
had a servant who was very reliable and to whom K. entrusted the carrying
out of all kinds of practical matters. Thus when he was to move house, he
would drive out in the morning, and in the evening he would go home to his
new apartment, where everything stood in perfect order --- even the
library had been arranged --- with the servant's help. Of this servant K.
once said, after speaking of his practical talent: ``he is in reality
\emph{my body}.''

\secnum{15.}

In the first period after Either/Or had come out, I had frequent
conversations with him about the book. On reading through the Διαψάλματα
I had at once recognized the anonymous author; I had at once found words
that had been spoken in earlier conversations with K. Naturally I never
tried, in my conversations with K., to get him to speak as the author; but
where something was obscure to me, I talked with him about it as with the
elder who had the richer insight. In this way I received many a fuller
elaboration, and he spoke about the book with a lack of reserve that he
would hardly have shown had I been indiscreet. One point on which he dwelt
several times was the poetic element in Part One. One day he explained
with great liveliness how in many places the motif for a poem was given,
but a motif that had surely been left unexecuted on purpose. Thus as an
example he singled out the depiction of woman's beauty and said:
% PRINTER: the Danish prints no opening quotation mark before "i hver af disse Sætninger" (the closing one after "Sonet." is printed); the English supplies the opening `` so that the marks balance.
``in each of these phrases: the little hand, the dainty foot, etc., there
lies a motif for a sonnet.'' --- Now and
% --- p. 349 ---
\opage{349}then he would hint at, as it were, the lines along which the
problems treated in Either/Or could be extended; but several of his hints
only became really intelligible to me when the subsequent writings came
out. --- He was not insensible to appreciative remarks about the book's
significance. I remarked one day that no book, since I first read Hegel's
Logic, had set my thoughts in motion as Either/Or had. This remark
evidently pleased him. We were just standing at the door of his lodgings
(on Nørregade) and were about to go our separate ways; as we parted, he
kindly gave me his hand, with the smile I knew so well.

\secnum{16.}

When I had published my piece against Martensen (``Nogle Bemærkninger om
Daaben''), K. talked with me about it. He was displeased with a review in
Berlingske Avis which had tried to strike a superior tone and, with an
assumed authority, meant as it were to crush the young author. I told him
that I had really only been amused by the review, since its forced
Christian piety took on a slightly comic tinge when one remembered that
the anonymous article appeared in a paper edited by a Jew. --- In my
little piece K. singled out one particular remark that he would have
wished left out; it was, I believe, where, with some mockery of
Martensen's attempt to assert the freedom of the newborn child, I say that
the only act of freedom that can be spoken of in the newborn child at
baptism is that the child screams when it is doused with water. This
remark, K. thought, looked flippant, and could easily be exploited by the
ill-disposed.

\secnum{17.}

K. mentioned Ludvig Feuerbach rather often in our conversations. It is on
the strength of oral remarks of his that I have said in Problemet om Tro
og Viden (p. 125) that there can be no doubt that it is Feuerbach whom K.
had
% --- p. 350 ---
\opage{350}in view in a quotation cited there. He acknowledged in
Feuerbach the clarity and sharpness with which he, precisely as one
offended, had grasped Christianity; but he once stressed in a conversation
that from Feuerbach's enthusiasm for nature, particularly from the pathos
with which in Wesen des Chr. he speaks of the strengthening power of
water, in the spiritual respect as well, he had received the impression
not merely of a strong sensuality in F., but also of an enfeeblement
through sensuality. He added:
% PRINTER: the Danish prints no opening quotation mark before "dette er en Mening" (the closing one after "Gisning." is printed); the English supplies the opening `` so that the marks balance.
``this is an opinion that I shall naturally never utter in print; but I
have not been able to ward off this impression, and there is something in
Feuerbach's peculiar passionate pathos, and in his antipathy to
Christianity, that becomes clearer to me through this conjecture.''
--- This remark about F. led him to speak of figures from antiquity in
whom a refined sensuality came to the fore. Thus he gave a very
imaginatively executed picture of Mæcenas: how he, who ``had no doubt
exhausted all the pleasures of Rome,'' sought oblivion in his
sleeplessness in the soft splashing of the fountain in his bedchamber and
in faint strains of distant music. Such pictures he could paint finely;
yet it was not the poetic depiction, but the psychological analysis, or
rather dissection, that interested him.

\secnum{18.}

On a longer walk in those years I once had occasion to observe with what
expression and what force K. could recite passages from a poetic work. He
was speaking of Aladdin, of the inspired daring in wishing that marks him,
of the contrast between this bold confidence of genius and the cautious
calculation of reflection. He now quoted Aladdin's words to the spirit,
where he commands him to build the splendid palace for the wedding night.
We were just then walking outside Charlottenborg. K. stopped as he
recited these lines, which he accompanied with
% --- p. 351 ---
\opage{351}a slight gesture of the arm. Color came into his cheeks, and
his gaze grew bright and warm, and his voice, though subdued, was that of
a ruler who knows that he has the power to command.

\secnum{19.}

Once we were talking about Hegel's presentation of the history of
philosophy. K. stressed how subjective its interpretation was on many
points, and said: ``geniuses really always lack the ability to grasp other
people's thoughts objectively; everywhere they find their own.'' --- Not
long afterwards, in a conversation with me, he happened to say of himself:
``I have never had this ability to grasp others objectively.'' --- At that
moment he hardly remembered the remark in our earlier conversation; I
remembered it, and smiled to myself as I drew the conclusion from his
remark about himself.

\secnum{20.}

At the time when Adler's mental derangement was beginning, K. told me a
couple of curious incidents about Mag. Adler. Adler came to K. one day
with a work he had published and talked with him at some length about the
religious authorship of them both. Adler gave K. to understand that he
regarded him as a kind of John the Baptist in relation to himself, who, as
the one who had received a direct revelation, was the real Messiah. I
still remember the smile with which K. told me that he had answered Adler
that he was perfectly content with the position Adler assigned him: he
found that it was a very respectable function to be a John the Baptist,
and he did not aspire to be the Messiah. During the same visit Adler read
aloud a large part of his work to K., some of it in his ordinary voice,
some in a peculiar whispering one. K. permitted himself a remark to the
effect that he could find no new
% --- p. 352 ---
\opage{352}revelation in Adler's
% PRINTER: the Danish prints "8krift" (a figure 8 for the initial S) for "Skrift"; rendered as the sense.
work, whereupon Adler said to him: ``Then I will come to you again this
evening and read the whole work to you in \emph{this} voice (the
whispering one), and then you will see that it will dawn on you.'' When he
told me the story, K. was much amused at this notion of Adler's that
varying the voice would make the work more significant. --- I once heard a
remark about Kierkegaard from Adler in my first year as a student. He
spoke of K.'s brilliant conversation, but put forward the opinion that K.
sometimes prepared himself for his conversations.

\secnum{21.}

As a younger man, and perhaps even after his dissertation, K. had thought
of the possibility of a post at the University. This thought, however,
soon receded. He once told me that Sibbern had urged him to apply for an
appointment as Docent in philosophy.
% PRINTER: the Danish prints "af han saa" for "at han saa"; rendered as the sense ("that ... then").
K. had replied that he would then at any rate have to stipulate a couple of
years' time for preparation. ``Oh! how can you think that they will
appoint you on those terms?'' said Sibbern then. ``Well, I could of
course also,'' retorted K., ``do as Rasmus Nielsen did, and get myself
appointed without being prepared.'' --- Sibbern grew cross and said:
% PRINTER: the Danish opens a quotation at "De skal" and never closes it; the English supplies the closing '' after "like that!" so that the marks balance.
``You always have to go pecking at Nielsen like that!'' --- It was a short
time after Nielsen had been appointed to the University.
---\,\dotfill\par
\noindent\dotfill\par
\noindent\dotfill

\secnum{22.}

In various of his writings K. has expressed himself on The Corsair's
attack on him. There he has viewed the matter essentially from an ethical
standpoint, and, as regards the significance the Corsair attacks had for
the general consciousness, he no doubt greatly overestimated
% --- p. 353 ---
\opage{353}it. K. did not have, in relation to a given phenomenon, a
sense of reality, if I may use this expression, that could form a
counterweight to his so enormously developed reflection. He could, as it
were, reflect a trifle up into something of world-historical importance.
So, no doubt, it went with him in the case of The Corsair. --- In a
conversation with me he spoke of the whole tendency of The Corsair and
judged its wit-mongering from a purely aesthetic standpoint. His chief
complaint against it was that it destroyed the public's sense for the
comic in its ideality by always attaching the laughter to particular
well-known persons who had once and for all been fixed for its readers as
comic figures, so that the reader laughed at what was said, not because
it really contained any comic substance, but because, without being really
comic, it was said of persons whom people had once and for all agreed to
find comic and to laugh at.

\secnum{23.}

Something like what happened to K. with The Corsair's attack happened to
him also with ``Søren Kirk'' in Hostrup's ``Genboerne.'' In his
posthumous journals K. has touched rather often on the fact that Hostrup
put him on the stage, and he has done so in a way that shows that he felt
affected by it. I played this role at the play's first performance, and I
remember having spoken with K. about it at the time without being able to
notice that he attached any importance to it; but in time he no doubt
reflected himself into a different mood. In the little role Hostrup took
motifs from S. K.'s writings, but in the role he was evidently thinking
chiefly of the kind of dialectic that was rampant among the young students
in those years, after ..... had given the impulse to a very superficial
preoccupation with philosophy. In this kind of dialectic K. was a master
in his younger years --- when he wanted to be. --- My conversation
% --- p. 354 ---
\opage{354}with K. about the role took place one evening when I was going
to a rehearsal at Hofteatret. On Højbroplads I met K., who stopped me and
asked where I was going. I told him I was going to the rehearsal of
Hostrup's play. ``Ah yes, you're to play me, of course!'' K. said then. I
answered him that I was not exactly going to play \emph{him}, but a
dialectician of our homegrown kind. We then parted without my detecting
that K. disliked my having taken on the little role. In any case he must
also have known my relation to him well enough to know that it could never
occur to me to put him in a comic light by copying him. For that I had
too much devotion to him, and --- for that I was much too poor an actor.

\secnum{24.}

In the early forties K. made a journey over to the west coast of Jutland,
to his father's native district. He told me a small characteristic
incident from this journey.
% "sin Fødeby": grammatically K.'s own native village, though the father's is meant; the English "his" keeps the looseness.
In his native village, for which his late father had founded a sizable
endowment for a school, he had been received with great deference, and
the schoolmaster in particular had been very busy in his official
capacity. ``I was really afraid they would put up triumphal arches for
me,'' he said. When he was to leave, he drove past the school. There
stood the schoolmaster, drawn up with all the schoolchildren, who were to
sing a song composed by the teacher in S. K.'s honor. The teacher, who was
to lead the singing, stood with a copy in his hand and was just about to
give the sign to begin when K.'s carriage stopped beside him. S. K. leaned
down, with his friendliest smile, toward the schoolmaster, took the copy
out of his hand, as if to read it through, and at the same moment gave
the coachman a sign to drive on. With that the whole arrangement fell to
pieces. The parish clerk, who did not know his poem by heart, could not
get the song started, the children stood silent
% --- p. 355 ---
\opage{355}and astonished, and S. K. rolled off down the highway, nodding
and greeting, and in his heart amused at the schoolmaster's
disappointment.

\secnum{25.}

On the day that Licentiat Hagen (later Professor of theology) had defended
his dissertation, I talked with K. about the dissertation (``Om
Ægteskabet''), against which I had been an opponent. K. was little pleased
with it; he found the whole treatment superficial, and
% Construal: "hvad der gives til en fyldigere Opfattelse af Pseudonymerne" read as what the pseudonyms offer toward a fuller conception (of marriage), not "a fuller conception of the pseudonyms".
he alluded to what is offered by the pseudonyms toward a fuller
conception. On this occasion too he stressed the Germans' more thorough
treatment of such subjects, and said with regard to Hagen's treatise:
``one really should see to it that we can cut a somewhat decent figure
when such books are read by the Germans.''

\secnum{26.}

While I was staying in Berlin in the summer of 1846, S. K. came there. It
was one of the little journeys he sometimes made when he really wanted to
work intensively and sought changed surroundings in order to put himself
in the mood. I met him quite unexpectedly at the restaurant where I ate,
and where he had gone chiefly to see fellow countrymen, since during his
first stay in Berlin this restaurant had been the gathering place of the
Danes. As a restaurant he did not think much of it; it confined itself, he
remarked, to ``the lowest concept of foddering.'' He met me with much
friendliness, invited me to dine with him the next day at the hotel, and
took long walks with me during the couple of days he stayed in Berlin. I
was much taken up with everything new that I came into contact with in
Berlin, and was tirelessly active. K. was amused at this activity, which
extended even to my attending the meetings of a democratic workers'
association. We talked a great deal together
% --- p. 356 ---
\opage{356}about conditions in Berlin, about conditions at home, about
my dissertation, etc. On our walks I noticed that there were certain spots
for which K. had acquired a fondness and which he liked to seek out. Thus
in the Thiergarten there was a pretty spot with a park surrounded by
flower beds. It became the goal of his walk, just as the flower beds at
the entrance to Frederiksberg Have had been in Copenhagen. Perhaps it
awakened memories in him of that garden and of what was connected with
it. --- Visiting him at the hotel, I noticed with what ingenious
calculation everything in his room was arranged so that he could work in
the right mood. The lighting, the connection between the rooms, the
arrangement of the furniture, everything was ordered according to a
definite plan. K. made very great demands on the service, but was also
very generous with his gratuities. Our little dinner for 2 persons was
served in his room. We drank a wine that he was fond of because of its
name: Liebfrauenmilch. This name had given him the impression of
something mild and light, and he drank the wine in good faith, trusting
this impression, without noticing that it is in fact one of the headier
Rhine wines. --- It did me good to meet K. in that foreign place. He
confirmed me in what I was then beginning to think about again: returning
to Denmark and trying to make my way there. --- From Berlin I once sent
K., who was always amused by the Berliners' Ueberschwänglichkeit, through
our cousin, Mrs.~T., an advertisement I had found in a Berlin advertising
gazette. It read, under the heading: Anfrage, as follows: ``Sollte es
nicht zweckmässig sein das auch die Klempner sich assoziirten, um mit dem
Zeitgeiste Schritt zu halten?'' --- It amused Kierkegaard, who
% "sine Pappenheimere": the proverbial allusion to Schiller (Wallensteins Tod, "Daran erkenn' ich meine Pappenheimer"), kept literally.
recognized his Pappenheimers.

% --- p. 357 ---
\secnum{27.}

\opage{357}After I had come home from my first Italian
journey (at the end of 1847), K. once came up to me
as I stood on Østergade looking at some plaster casts
of antique sculptures that stood in the window of the plaster-caster
Barsugli. When we came to talk of the art of sculpture, he remarked
that he had not really had an eye for
it, but that it was now beginning to attract him by
its repose. On the whole I do not believe that K.'s interest in
and feeling for the visual arts were strongly developed; his
interest in the beautiful drew its nourishment overwhelmingly
from poetry.

\secnum{28.}

Shortly after the rebellion had broken out in 1848, I talked
with K. about the plan I then had of going along
as a volunteer. Besides the general
motive that lay in the circumstances, I stressed that I personally
needed to gain experience, above all experience of myself. K.
granted that this way of learning through reality was
necessary for many, but intimated that through
imagination and reflection he took in, ideally, a still richer
content. And it was evidently the case that
he could make do with a \emph{motif} from reality, which
his reflection and his imagination worked up, guided as they
were by the experience he had gathered through the ceaseless
labor of his own inner life.

\secnum{29.}

Some years later I talked with K. about a similar
subject. At that time I felt dissatisfied at not
having a definite occupation, and the definite duties
that go with one; the complete freedom I
enjoyed in my life of private study, holding no office, was beginning to become
a burden to me. I said so in a conversation with K.:
``I miss a definite practical occupation to such a degree
% --- p. 358 ---
\opage{358}that, if no other is to be found, I could be tempted to set
myself up as a tavern-keeper.''---K. smiled and said:
``Then you shall have my custom.'' In that connection
he told me how, some years earlier, he had received
a visit from a German scholar who had presumably been sent to him
by someone or other to look at him as one
of our curiosities. He had received the German
very politely, but had assured him that there must be
some misunderstanding behind the visit. ``My brother, the Doctor,''
he had said, ``is an extraordinarily learned man
whom it would surely interest you to get to know,
but I am a beer-seller.''

\secnum{30.}

Dr. P. K. once gave---I do not remember the time exactly,
but it must have been in the fifties---during a chance
stay down here, a series of lectures at the
University. They were held in the Solennitetssal before a large
and very mixed audience. Søren K. took
a rather ironic attitude toward these lectures, and it was with
a certain malice that he told me a fragment of a
conversation he had overheard one evening outside the University,
just as the lecture was beginning. A crowd of
gentlemen and still more ladies came streaming up the University steps.
The coachman of a carriage waiting outside was
addressed by a passing fellow with the question
of what all these people were going in there for? ``Oh! they'll be going
to a dance, I suppose!'' answered the coachman, to S. K.'s great delight.
The lectures were in the Doctor's manner, rather
forcedly brilliant, and in a prophetic-apocalyptic style. I
was able to tell S. K. an anecdote concerning the lectures,
about our old uncle in Købmagergade.
He attended the lectures regularly, although they lay
far above his horizon and although he was very
hard of hearing. He felt exceedingly proud to see his
``cousin's son on Gammeltorv'' lecturing to so numerous a
% --- p. 359 ---
\opage{359}crowd. One evening, when he had come home from the lectures,
he walked up and down the floor of his parlor,
repeating without pause: ``No! he really is a matchless man,
that Peter Kierkegaard! the way he can talk, the way
he can talk.'' His daughter ventured the
question: ``Now, could you really follow his lecture?
% The Jutlander says "en Ord" for standard "et Ord" (dialect gender); the
% English cannot carry the dialect -- hence "den gamle Jyde" just after.
could you really hear what he said?''---``Not \emph{one}
word! not \emph{one} word!'' said the old Jutlander guilelessly,
but his enthusiasm for the matchless speaker was just as great.

\secnum{31.}

S. K. once told me of a curious situation involving this old uncle.
The old uncle had a passion
for making verses, equally ghastly in content
and form. One day he came up to see S. K. and, after a
little preamble, came out with a pile of his
verses, which he asked his nephew to read aloud. They then took
their places side by side on the sofa. The old man sat
leaning back and put on his spectacles so as to
follow along during the reading, most likely in order to
make sure that nothing was skipped. Søren K. sat
somewhat bent forward, with the papers, and presumably chose
this position so that the old man should not observe a
treacherous smile on his face. He read the poems
aloud from one end to the other in a raised voice and with great
pathos. The old man was utterly enraptured; the tears ran
down his cheeks from emotion at the beautiful verses,
and he left S. K. with the warmest thanks.---
Old Uncle's verses I have often heard; they were altogether
remarkable, and had in particular the wonderful property
of being usable, or at any rate
of being used, with very slight variations, on the most heterogeneous occasions.
Thus the old man had once, on the occasion of a
granddaughter's engagement, composed a verse. It was then
used not only at another engagement, where the personal
circumstances of those concerned were, after all, very different;
% --- p. 360 ---
\opage{360}but with a quite trifling adjustment it was also
to have been used on the occasion of my appointment
to the University. That it was not was owing only
to the old man's not having got it rearranged on the
evening when the appointment, a little earlier than we had expected,
was announced by Berlingske, and we happened to be at a party at
the old man's.---Sometimes the verses were sung by the company
at his house at his urgent request. A
tune would then be chosen which by a Procrustean method
could more or less be applied to them; but it was irresistibly comic
to hear how now a large number of syllables
had to be taken, as it were, in one mouthful, now a single
syllable drawn out according to the Holbergian prescription.
But to the old man's ears it always sounded like the loveliest
harmonies; his face beamed with rapture
at it.

\secnum{32.}

Old Kierkegaard sometimes used, in speech and in
writing, though chiefly in writing, peculiar orders of words. Thus,
when in his letters to his native town he
spoke of S. K. and his brother, he used always to write: ``my
cousin's late son Peter K., or: my cousin's late
son S. K.'' (instead of: my late cousin's son). It
looked especially comic in one context: when he reported
that his cousin's late son, Peter Kierkegaard,
had taken himself a new wife. I once told S. K.
of this way of putting things on the uncle's part. It amused him,
but for his own part he read into it a meaning of which the
old man was unconscious (the old man was, after all, inspired, as a
poet!): that he was in reality one of the departed.
---On the same occasion he came to speak of
the uncle's great age. I said that in a certain sense
one could not call him old, if, that is, one
measured the length of life by its content, since the old man
in reality only ``lived'' the smallest part of the
% --- p. 361 ---
\opage{361}time he existed, and for the most part merely vegetated and slept.
I added jokingly that S. K. was really the oldest
man I had known. He smiled and made no
objection to this way of reckoning.

\secnum{33.}

K. once told me---it occurs to me as I
mention his age---that as a young man he had for many years
lived in the firm conviction that he would
die when he reached the age of 33. (Was it the age of Jesus
that was also to be the norm for the imitator of Jesus?)---
% "Denne Tro": ordinary belief here, not faith in the religious sense;
% rendered "belief" (cf. "at tro dette", "to believe this", below).
This belief was so firm in him that, when he passed
that age, he even had the parish register looked up in order
to make sure that it really had been passed; so
hard was it for him to believe this.

\secnum{34.}

When I first lectured at the University,
in the winter semester of 1849--50, S. K. talked with me about it
fairly often, about the subject of the lectures (Greek
philosophy) and about the plan for treating it. He
told me that he would come round some time and hear
them, but that he would not name any particular hour.
He did not come, however; but through that constant prospect
of possibly having him as a listener he gave me an impulse
to treat my subject with care, and
I rather think that that is what he intended. He also
knew that, if he really were present, I
would easily feel embarrassed by his presence, and that may
have been one reason for him to stay away.

\secnum{35.}

In connection with my lectures he once talked
with me about the Aristotelian metaphysics' definition
% "den aristoteleske Metafysiks Bestemmelse": so the Danish, though the
% definition stands in the Nicomachean Ethics (named below); rendered as written.
of virtue as μεσότης. He had rightly seen that this
definition did not apply to virtue in general, but
% --- p. 362 ---
\opage{362}only to ethical virtue---something on which people at that time were by no means
clear. He told me that, in connection
with this point and some others in the Nicomachean
Ethics, he had sought information from Madvig and Sibbern.
Madvig had told him that it was so long since
he had read this work that he no longer had it
% "et Jærn" (lit. "an iron"): colloquial for a person of formidable, tireless
% capacity; rendered by sense.
present to mind. ``And Madvig is otherwise a formidable man,'' S. K.
added. Sibbern, on the other hand, had at once been able to give him the
information he sought. S. K. thought highly of S., even if
he was not blind to his weaknesses. Among
these he once singled out Sibbern's complete
lack of irony and---in psychological respects---his
lack of an eye for the masked passions,
for the doubling by which one passion
takes on the form of another. Because of this, in his opinion,
Sibbern had often been deceived at a time when many, especially
ladies, turned to him to consult him as
a kind of psychological counselor of souls.

\secnum{36.}

Poul Møller was someone S. K. often came to speak of, always
with the most heartfelt affection. It was Poul Møller's
personality, far more than his writings, that had
made an impression on him, and he regretted that a time would
soon come when, once the preserved memory
of his personality had vanished and judgment
of him rested only on his writings, people would no longer
be able to understand his significance. He once told me
an amusing little anecdote of P. Møller. He was to
oppose ex officio at a disputation and had noted down
his remarks on various loose sheets which had been
laid inside the dissertation. Each separate objection he
introduced with a ``graviter vituperandum est''; but as soon as
the Præses had given an answer to the objection, he would say
very good-naturedly: Concedo, and pass on to the next
objection. After a rather short opposition he closed
% --- p. 363 ---
\opage{363}with a lively expression of regret that his time did not
permit him to continue this interesting conversation, and
as he then withdrew and passed S. K., who stood
among the audience, he said to him in an undertone: ``Shall we
go down to Pleisch's, then?''---that was the confectioner's where
he used to go.---During the opposition it had happened
to him that all his loose papers had at once fallen
out of the book and fluttered about on the floor, and it
had contributed not a little to raising the mood of the audience
to see that large figure crawling about to
gather up the scattered sheets.

\secnum{37.}

A man with whom in those years I regularly saw K.
walking was Professor Kolderup-Rosenvinge. It was
chiefly a shared interest in the aesthetic that brought
them into contact. K.-R. was a cultivated man who occupied
himself with the literatures of southern Europe and had
translated a play by Calderon. For the rest, he was by then
already rather dulled, and in many respects
extremely narrow-minded. When I once expressed to S. K. my
astonishment that he could find it interesting to talk
with Rosenvinge, he stressed his general cultivation.
In general he set great store by people of
the older generation who had preserved the humane interests
of an earlier time and that refinement of manner which a younger
generation so much lacked.

\secnum{38.}

With the late Pastor Spang K. had for a time (in the years
1844--45?) associated a great deal, and had almost daily
called for Spang at his home for a longish evening walk.
(Several interesting letters to Spang from S. K. are
in the possession of Spang's daughter, Mrs.~Rump). After Spang had died,
K. once told me how his widow had been
overwhelmed with grief, and how he had, by his
% --- p. 364 ---
\opage{364}conversation, helped to bring her mind back into balance.
On that occasion he used the expression: ``I
know how one should grieve, and how one should
comfort in grief.'' And indeed: he understood it as few do.
He comforted, not by covering up the grief, but by
first bringing it properly to consciousness, clearing it up completely,
and then reminding one that it was a \emph{duty} to grieve,
but for that very reason also a duty not to let oneself be crushed
by grief, but in grief to preserve the strength for one's
work, indeed to find in grief the incentive to do that work
still more fully.---I have myself, in my
relationship with S. K., more than once had the experience of
how he knew how to lift one up when one was bowed down,
to comfort when one grieved, and that without one's
needing to speak of what was oppressing or saddening one.
Thus I remember how once, in the autumn of 1850, when,
because I had not been able to find lodgings in time,
I had to stay for a short while at a cheerless inn,
% Danish prints "Iudtrykket" (turned n) for "Indtrykket"; rendered as the sense.
and under the impression of my surroundings despondent thoughts
had got the upper hand, I then one evening, as I was walking dejectedly
in the street, met K., who began a conversation
with me. Without my needing to say a word about
myself, he saw with his sharp eye that I needed
to be brought out of an oppressive mood, and he knew how,
by his conversation, without its seeming for a moment to be directly
aimed at that, so to free my mind that I left
him glad and in good heart, and for a long time was freed from
the power of melancholy.

\secnum{39.}

Much of what S. K. worked out in his writings
I remember him touching on in conversation with me while
his thought was at work on it, and precisely because
it was still in the making, it had a life and an attraction
even greater than what it could have when fully formed.
Thus I remember that once, while he was walking along
% --- p. 365 ---
\opage{365}with Christens and me, he developed the theme he so often
has treated in his upbuilding writings, that ``the \emph{whole} of life
is the time of trial.'' He developed it in this form: that
for the believer there was an ``At last'' which seemed to
be so near that it could be grasped, but constantly drew
back, until the end of life finally led to grasping
it. Only in death could it be said: ``At last!''---
Sometimes in his conversations S. K. expressed his thought in
a more pointed, more paradoxical form than in his writings.
It amused him to formulate it so that it formed
the opposite of something commonly accepted. Thus, in opposition
to the maxim: experience makes one wise, he used
to put forward the theory: ``Experience makes one mad!'' and he
could carry out the proof with much humor by pointing
to all the countless contradictions of concrete experience.

\secnum{40.}

The different direction of our spiritual lives and our
fundamentally different relation to what is Christian came now and then
% Danish prints a comma after "sammen" where the sentence ends (capital "Engang"
% follows); rendered as the sense, with a full stop.
distinctly to light when we talked together. Once, probably
in 1851, he asked me what I was occupied
with at the moment. I answered him that I was reading the
N. T. That seemed to please him, but it did not please
him when I added that I was reading it chiefly with
the aim of finding, in what has been handed down, what is primitively
Christian, and of following the successive development of
Christian dogmatic ideas. For him such
an investigation was something almost offensive. Although
he draws the dogmas into his writings so little, and everywhere
keeps to what is central in Christianity, he
could not, for all that, place himself in such a relation to Scripture
that it became the object of a critical investigation. He
could not do so, partly because Scripture stood for him as a
unity, as the unified expression of what is Christian, into which
he would bring no distinction; partly because, by
% --- p. 366 ---
\opage{366}being made the object of a scientific investigation,
it became an object of knowledge instead of an object of faith.

\secnum{41.}

% Danish prints "Forboldet" (b for h) for "Forholdet"; rendered as the sense.
My divergence from S. K. in relation to what is Christian
I did not hide from him; but we did not dispute
about it. There was so extraordinarily much in which we
could be in sympathy, and I preferred to let myself be taught
by him about much that was of great importance to me,
rather than carry on a discussion in which, by his superiority,
he could easily have overwhelmed me, but not
easily convinced me. In our conversations, too, his attitude toward me
was essentially communicative, and I saw in this a
proof that he bore me goodwill. He had
two ways of conversing: the one essentially communicative,
awakening, encouraging; the other questioning,
ironizing, and, through his dialectic, confounding.
The latter he never used against me.
Once in a while, when I grew heated, he might throw in a
joking, teasing word. Thus I was once speaking
very heatedly about how no positive religion could be
tolerant, precisely because, by its claim to be
revealed religion, it had to lay claim to being
the only true one, and therefore had to posit the others as
untrue. Something universally religious, a ``religion in general,''
must therefore, from the standpoint of positive religion, be
an absurdity. As I was developing this with heat, I happened
to repeat the expression: a religion in general,
and made my own the proposition that a religion (i.e.
% Pun the English cannot carry: "overhovedet" ("in general", "at all") is
% literally "over the head"; K. answers "over Hovedet", "over the head".
\emph{positive} religion) in general was an absurdity.---``Yes, and
a chamber pot over the head,'' said K., and thereby damped
my ardor. That is, however, the only piece of teasing of
that kind I remember from him.---If he never tried
to counter my divergent views directly,
I can, given the whole friendly sympathy he always showed
me, and which I gratefully remember, only seek the reason
% --- p. 367 ---
% p. 367 opens mid-sentence (Danish "kun søge Grunden / til det deri, at ...").
\opage{367}for it in this, that he knew that I was occupying myself honestly and
earnestly with the cause that was for him the highest, and that I was so
familiar with what he had written about it, and with his whole way of
thinking, that I had sufficient premises for a conclusion, and would not be
brought nearer to it by another, but in any case only by myself.

\secnum{42.}

Concerning the general effect of S. K.'s writings I once said something to
him that seemed to make some impression on him. He was speaking of how, as an
author, he had had the peculiar fate of publishing one work after another
without calling forth any real criticism, whether hostile or approving. He
then told me that his writings had nevertheless found a circle of readers
from the beginning, and, it seemed, a fixed circle, since the number of
copies sold straight away was roughly the same (c. 70 for the writings that
followed Either/Or). I expressed my conviction that a great many people read
his writings with real seriousness, and that the silence of our press, which
I explained chiefly by the critics' feeling that the writings lay beyond the
powers of their criticism, was by no means a sign that he was meeting with
indifference in a wider public, on which I believed that he was, precisely,
already then having an effect on a large scale. But I added that the effect,
as regards the relation to Christianity, had in my opinion been as much
negative as positive. All those on whom the writings had had an effect, in
the one direction or the other, were surely agreed that only in his
conception did Christianity have a real distinctiveness and inspire
reverence. But his definition of Christianity was such that it must,
precisely, have the power to
% --- p. 368 ---
\opage{368}drive many away from it, as was the case with me, and must have it
precisely through the separation it set between Christianity and nature and
concrete human life. --- We were, when I said these words to K., just in the
act of taking leave of each other; he grew serious and silently gave me his
hand, and parted from me, but with a friendly expression on his face.

\secnum{43.}

S. K. spoke without reserve of the significance his writings could have for
certain particular persons. His old uncle, the wholesale merchant M.
Kierkegaard, had a son, a couple of years younger than S. K., who was
disabled, paralyzed down the whole of one side, and wholly a cripple;
spiritually, on the other hand, well gifted. He read his cousin's writings
with the greatest interest, visited him now and then at his home, and drew
from these visits much spiritual uplift. I once talked about him with S. K.
and told him how great an impression one of his works, the confessional
discourse in ``Upbuilding Discourses in Various Spirits'', had made on him. In
it there is a passage about a person who by his bodily weakness is prevented
from carrying on any work in the outward world, and it is depicted beautifully
and upliftingly how such a person too has the universally human ethical task
undiminished, and what particular form the task of life takes for him. S. K.
said: ``Yes! for him this book is a blessing.'' And so it truly was. It had
the power to give the sorely tried man strength to overcome the thought that
his existence was useless and wasted, and to give him the consciousness of
being essentially the equal of those more happily endowed by nature. Precisely
this consciousness S. K. also knew how to call to life in him by his
conversations; that is why he came away from them so strengthened.

% --- p. 369 ---
\secnum{\opage{369}44.}

A conversation that I have often remembered since, I had with S. K. some time
before his ``Practice in Christianity'' came out. We were walking together
down by the Lakes, and he urged me not to put it off too long, beginning an
authorship. He added jokingly: ``Now, you know, there is a place vacant in our
literature; \emph{now I am finished}.'' This remark, which I passed over
lightly at the time, came back vividly to my memory later, when the last
section of S. K.'s authorship lay completed. When he spoke to me in this way,
he evidently still cherished the hope of being able to avoid the last battle
with the established order. Had that battle not been forced upon him, his
authorship would indeed have been --- essentially ended, concluded, when that
conversation took place. In reality all the stages that were to be traversed
had by then been traversed; his writings formed a wonderfully complete whole.
But the next stage still bore within it a possibility of a redoubling, and
this possibility became actuality when an attempt was made to reduce his
conception of Christianity to an eccentric exaggeration, to set up mediocrity
and spiritlessness as what is properly normal.

\secnum{45.}

In the same conversation S. K. came to single out one aspect of his
authorship, which he spoke of with the objectivity peculiar to him: it was
the significance of his writings for Danish prose. He said: ``In this century
our literature has displayed an almost abnormal wealth in the development of
its poetry; our prose literature, on the other hand, has lagged behind; we
lacked a prose that bore the stamp of the artistic. This lack I have filled,
and that significance my writings will also retain for our literature.'' In
this I believe he is right: his prose is art, not without its shadows, but as
a whole the clearest,
% --- p. 370 ---
\opage{370}most supple, most comprehensive and fullest expression for thought,
and no less for feeling and passion, that our literature knows. It has often
reminded me of Plato, the author with whose prose I would most readily
compare Kierkegaard's, and who has surely also given him the model, which he
has freely made over in his own way.

\secnum{46.}

In his writings S. K. has in several places spoken of Grundtvig and
Grundtvigianism, and almost everywhere seen it from the comic side. In a
similar way he often expressed himself about it in his conversations. The
sharp contrast between his conception of the religious and the Grundtvigian
one constantly came out, and the Grundtvigian childish immediacy was cast in
a comic light by the contrast. Only about one period of Grundtvig's life did
S. K. once speak to me with interest and appreciation: it was the time when
Gr. came out against Clausen. But what appealed to K. in his conduct then was
evidently not the specifically Grundtvigian, least of all the ``matchless
discovery'', which always seemed to K. a matchless piece of foolishness, and
an unoriginal piece of foolishness at that, but what forms the
opposite of rationalism: the religious stirring of \emph{life}. S. K.'s
father, whose religious orientation was closest to that of the old Pietism,
sympathized with Grundtvig and Lindberg, surely from this point of view, and
S. K.'s sympathy undoubtedly had the same motive and the same limits as his
father's.

\secnum{47.}

A little trait of S. K.'s comes to mind as I think of Grundtvig. The
patroness of Grundtvigianism, the Queen Dowager Caroline Amalie, also thought
highly of S. K., at least before his attack on Mynster. S. K. was evidently
not insensible to goodwill shown
% --- p. 371 ---
\opage{371}him by a queen. One day I was walking with him on Kærlighedsstien,
when he suddenly burst out: ``Well, devil take it! now I shan't get away.'' I
looked at him in surprise, for otherwise he never used words resembling an
oath, and looked ahead along the path, where a long way off I caught sight of
a footman in red livery, and in front of him two ladies. One of these was the
Queen Dowager, who, when we met her, very graciously stopped S. K. and talked
with him. When he saw her, we were so far from her that we could easily have
turned back without having been noticed. That his outburst was not altogether
sincerely meant was the impression I got from this: ``well, devil take it!''
was not like him, and it came out somewhat forced, as when one wants by a
strong expression to hide that one is really quite pleased that something is
happening. S. K.'s old, traditional reverence for the royal power has not
been without significance in his relation to the Queen Dowager.

\secnum{48.}

About Bishop Mynster I had, in the spring of 1852, a conversation with S. K.
that I often could not help thinking of later, when K. came out against him.
One afternoon I was walking on the Ramparts, on the promenade that runs along
the top of the parapet, and there I passed Bishop Mynster. A moment later I met
S. K., who was walking below on the Ramparts and called to me. He had seen
Mynster walking ahead of him, and so our conversation turned to him. K. first
spoke of his imposing and venerable appearance, and then of his shrewd and
polished manner, and how as a man of the world he was superior to most of the
younger men, who for a time had believed that with their philosophical
apparatus they could treat Mynster as a subordinate figure, but soon, on
personal contact with him, had been overwhelmed by him and now subordinated
themselves to him. He named in particular
. . . . . as one who by his awkwardness and lack of polish gave Mynster so
high a degree of
% --- p. 372 ---
\opage{372}superiority in their personal dealings. But he stressed very
strongly that Mynster's forte was precisely his worldly wisdom, and that it
was this that determined his conduct. Of his own relation to Mynster K. spoke
ironically. He told me with a smile that he had once asked Mynster's
permission, whenever there were problems that he could not himself bring to
clarity, to seek guidance then from His Right Reverend Excellency! and that
Mynster had kindly given him this permission. The smile alone with which he
told me this showed me clearly how little, in reality, he thought of Mynster
at that time, and the whole manner in which he characterized Mynster as a man
of the world agreed perfectly with his later pronouncements. There was a time
when he had thought very highly of Mynster, and he had come to this
essentially through filial piety toward his father, who valued Mynster
greatly. This dutiful regard for his father's judgment I have seen in several
cases determining K.'s view. Thus I remember that he once spoke with much
goodwill of a man who in no respect deserved esteem or could by his person
have aroused S. K.'s interest; but his father felt great attachment and
gratitude toward this man, who had once, by a shrewd piece of advice,
prevented his financial affairs, which were endangered in the confused
monetary period after the war with England, from falling into disorder.

\secnum{49.}

When S. K. began his polemic against the established order, and perhaps some
time before that, he stopped taking part in the public worship of the Church,
at which he had previously been a very regular attendant. The late Dr.
Frederik Beck once told me that in that period S. K. liked to come into the
Athenæum on Sunday mornings during church time; which Beck, with the relish
for barbed remarks peculiar to him, interpreted
% --- p. 373 ---
\opage{373}to mean that S. K. wanted to make people notice that he did not go
to church. --- Once, many years earlier, S. K. had asked me whether I went to
church. I said no, and told him that I did not, partly because I felt a
stranger there, partly because the service in the Copenhagen churches went
on amid so many disturbances that I could in no case keep a devout frame of
mind, but was more likely to be driven to take offense. He then said that it
was a \emph{duty} not to let oneself be disturbed by what went on around one;
one \emph{should} be upbuilt when one was in church.

\secnum{50.}

The last time I spoke with S. K. was in the summer of 1855, when I had not
seen him for an unusually long time. It was one evening as I was walking from
Højbroplads up toward Vimmelskaftet. K. came up to me and began a
conversation. After we had talked for a moment about my brief part in
political life (in the winter of 1854---55 I had been a member of the
Rigsdag), we soon came to speak of his polemic against the established order
of the Church. I thanked him warmly for his stand, and told him how many I
had found who sympathized with him wholeheartedly. He spoke of the situation
he had brought about with the greatest clarity and calm, and it astonished me
that in the midst of the violent struggle, which cut so deep into his life
and laid claim to his last strength, he could preserve not only his wonted
composure and confidence, but could still keep his sense of fun. As we were
nearing his home --- at that time he lived in Klædeboderne --- he said in a
joking tone: ``Now I'll go home and go to bed, and then I'll say to the world
what that Kammerraad said.'' I asked him what, then, this Kammerraad had said.
He then told me that there had once lived an eccentric there in the street,
who was a Kammerraad by title, and who regularly every evening
% --- p. 374 ---
\opage{374}used to hold the following conversation with the watchman. When
the watchman had called ten, the Kammerraad would open the window and ask:
``Watchman, what time is it now?'' --- ``It is ten o'clock, Mr. Kammerraad.''
% The Danish never closes the quotation that begins "Godt!" (and prints its opening
% mark as two commas). The English closes it after the break-off dash, before the
% parenthesis, where the sense requires, so that its quotation marks balance.
--- ``Good! then I think I'll go to bed. If anyone should ask for me now,
watchman, you may tell them to ---''
% Götz von Berlichingen's reply (in Goethe's play) is the proverbial "kiss my arse"; the Danish leaves it unspoken.
(Götz v. Berlichingen's answer.) --- He told this story with his old humor,
and I bade him Good night! without suspecting that I had that evening spoken
with him for the last time.

\begin{center}\small (Written between 27 Dec. 1871 and 10 January 1872.)\end{center}

\begin{flushright}\textbf{H. Brøchner.}\end{flushright}

\end{document}
